\subsection{Recovering trust for contactless PIV Auto}
\label{ex:sm-recovery}

This example uses a credential that supports SM and requires a pairing code
for VCI. PIV Auto is configured for contactless authentication with key
\code{0x9A}, but the PD's trust store is empty. The ACU already holds the roots
and validation policy authorized for the deployment. Figure~\ref{fig:sm-recovery}
shows the recovery sequence.

\begin{figure}[htbp]
\centering
\begin{tikzpicture}[x=1cm,y=1cm,font=\sffamily\small,
  recoverybox/.style={draw=OSDPBlue,rounded corners=3pt,fill=OSDPLightGray,
    text width=12.5cm,align=left,inner sep=9pt},
  flow/.style={->,draw=OSDPBlue,thick}]
\node[recoverybox] (failure) at (0,0) {
  \textbf{1. PD reports why PIV Auto stopped}\\[3pt]
  This PIN-based contactless 0x9A workflow requires SM, VCI, and a verified PIN.
  An empty trust store produces PIVSTATUS Result 0x02, Detail 0x01.
  The PD ends this attempt.};
\node[recoverybox,anchor=north] (inspect) at ([yshift=-0.45cm]failure.south) {
  \textbf{2. ACU obtains the diagnostic and signer information}\\[3pt]
  CARDSTATUS Type 0x02 requests CARDSTATUSR Type 0x02 fragments containing
  the setup outcome and any received CVC's hash.
  PIVGETDATA for 0x5FC122 reads the signer object from the credential.
  Recovery stops if that read fails. An empty store means setup was skipped, so the
  diagnostic reports that no CVC hash is available.};
\node[recoverybox,anchor=north] (trust) at ([yshift=-0.45cm]inspect.south) {
  \textbf{3. ACU validates and loads the trust anchor}\\[3pt]
  The ACU validates the signer certificate against an authorized root and
  deployment policy. After validation, PIVVCILOADTA loads the anchor.
  The ACU reapplies PIVMODE to start a new automatic attempt.};
\node[recoverybox,anchor=north] (pair) at ([yshift=-0.45cm]trust.south) {
  \textbf{4. PD reports pairing required; ACU supplies the code}\\[3pt]
  SM succeeds, but this credential requires pairing. PIV Auto reports
  Result 0x05, Detail 0x01 and ends the attempt. The ACU sends PIVXMITPAIRING
  and waits for status confirming SM and VCI.};
\node[recoverybox,anchor=north] (auth) at ([yshift=-0.45cm]pair.south) {
  \textbf{5. ACU completes authentication in commanded mode}\\[3pt]
  PIVSPE Mode 0x01 collects and verifies the PIN. The ACU obtains the 0x9A
  certificate with PIVGETDATA and sends a fresh challenge with CRAUTH.
  The ACU validates the certificate and response before deciding access.
  The existing SM and PIN state remain in use.};
\draw[flow] (failure.south) -- (inspect.north);
\draw[flow] (inspect.south) -- (trust.north);
\draw[flow] (trust.south) -- (pair.north);
\draw[flow] (pair.south) -- (auth.north);
\end{tikzpicture}
\caption{Recovery from an empty PD trust store for contactless PIV authentication
with key 0x9A. Every step depends on a successful result from the preceding step.
Unsupported SM or VCI requires a different ACU-selected workflow.}
\label{fig:sm-recovery}
\end{figure}


The first PIV Auto attempt ends with Result \code{0x02}, Detail \code{0x01}
(trust store empty). The ACU requests the SM diagnostic record, then reads the
Secure Messaging Certificate Signer object with PIVGETDATA, Object Identifier
\code{0x5FC122}, Tag Length zero, Offset zero, and Requested Length zero.
When the read succeeds, the reply contains the complete object, including
signer certificate tag \code{0x70} and any intermediate CVC in tag \code{0x7F21}.
The PD reads the credential through the contactless interface. If the object
is absent or the credential denies access, the PD reports the corresponding
PIVERROR. The ACU then selects another supported workflow or ends the
transaction. Recovery relies on the credential returning the object.

The ACU validates the signer certificate's chain to an authorized root,
validity period, revocation status, and PIV or PIV-I content-signing extended
key usage. The ACU checks any intermediate CVC's signature and linkage as
specified under PIVVCILOADTA. The anchor record contains the validated signer's
public key and the IIN derived from that certificate's subjectKeyIdentifier.
The ACU loads an anchor only after these checks succeed.
An IIN match alone provides no authorization to trust a certificate.

After the load succeeds, the ACU reapplies the PIVMODE configuration. The PD
starts a new attempt, establishes SM, and reports Result \code{0x05}, Detail
\code{0x01} (pairing code required). The ACU sends PIVXMITPAIRING and waits for
CARDSTATUSR with security status \code{0x01} (SM and VCI established).
The ACU then sends PIVSPE Mode \code{0x01}, with Auto Execute clear, to collect
and verify the PIN. After successful verification, the ACU reads the PIV
Authentication certificate (Object Identifier \code{0x5FC105}) and uses CRAUTH
with key \code{0x9A} and a fresh challenge. Certificate and challenge-response
validation remain the ACU's responsibility.

The ACU waits for a successful result at each step before issuing the next
command. An error ends this recovery sequence. These are commanded operations
after the automatic attempt has ended. The PD
returns each command's defined reply; the sequence generates no second PIV Auto
result. PIVMODE remains configured for the next credential presentation.
Reapplying PIVMODE after pairing or PIN verification would clear the state
needed for the commanded authentication.

If SM or VCI is unsupported, the ACU chooses a supported interface or
authentication method explicitly, or rejects the transaction. If a signature
or cryptogram check fails, the PD blocks credential operations and retains
the diagnostic record. The ACU investigates that failure before requesting
a new attempt. Trust recovery is appropriate for missing authorized trust
material; failed cryptographic checks are security failures.

This is a control-flow example. The preceding captured SM/VCI packet examples
cover pairing and PIN verification. The companion
\href{run:worked-examples/sm-vci/trust-recovery.md}{trust-recovery example}
records the command sequence and validation decisions.
